\documentclass[a4paper,12pt]{article}
\usepackage[utf8]{inputenc}
\usepackage[polish]{babel}
\usepackage[T1]{fontenc}
\usepackage{amsmath}
\usepackage{graphicx}
\usepackage{geometry}
\usepackage{float}
\usepackage{caption}
\usepackage{listings}
\usepackage{xcolor}
\usepackage{subfigure}
\usepackage{booktabs}
\usepackage{array}

\geometry{margin=2.5cm}

% Styl dla kodu MATLAB
\lstdefinestyle{matlab}{
    language=Matlab,
    basicstyle=\ttfamily\small,
    keywordstyle=\color{blue}\bfseries,
    commentstyle=\color{gray}\itshape,
    stringstyle=\color{orange},
    numbers=left,
    numberstyle=\tiny\color{gray},
    numbersep=8pt,
    frame=single,
    breaklines=true,
    backgroundcolor=\color{gray!8},
    captionpos=b
}

\title{
\Large Zaawansowane metody przetwarzania i analizy sygnałów \\[6pt]
\textbf{Analiza częstotliwościowa HRV i klasyfikacja grup wiekowych na podstawie bazy Fantasia Database}
}

\author{Piotr Sawicki 319003, Michał Odziemkowski 311267}
\date{Czerwiec 2026}

\begin{document}

\maketitle

% =========================================================

\section{Cel projektu}

Celem projektu była analiza częstotliwościowa zmienności rytmu serca (ang. \textit{Heart Rate Variability}, HRV) na podstawie zapisów elektrokardiograficznych pochodzących z bazy danych Fantasia Database. Analiza obejmowała wyznaczenie parametrów widmowych LF oraz HF, opisujących aktywność autonomicznego układu nerwowego.

Drugim celem pracy było opracowanie modelu klasyfikacyjnego umożliwiającego automatyczne rozróżnianie osób młodych i starszych na podstawie cech wyznaczonych z sygnału HRV. W tym celu zaimplementowano i porównano dwa klasyfikatory: Support Vector Machine (SVM) z jądrem RBF oraz Linear Discriminant Analysis (LDA).

W celu zwiększenia liczby próbek uczących zastosowano segmentację 120-minutowych nagrań na niezależne okna 5-minutowe, zgodnie z zaleceniami dotyczącymi krótkoterminowej analizy HRV.

% =========================================================

\section{Metodyka badań}

W celu realizacji założeń projektu opracowano kompletny łańcuch przetwarzania danych obejmujący ekstrakcję sygnału HRV, wyznaczenie parametrów częstotliwościowych oraz klasyfikację grup wiekowych. Analiza została przeprowadzona w środowisku MATLAB z wykorzystaniem biblioteki WFDB Toolbox umożliwiającej odczyt danych z platformy PhysioNet.

\subsection{Baza danych}

W badaniach wykorzystano bazę \textit{Fantasia Database} udostępnioną przez platformę PhysioNet \cite{fantasia,physionet}. Zbiór zawiera 40 długoterminowych zapisów sygnału EKG zarejestrowanych u zdrowych ochotników podzielonych na dwie grupy wiekowe:

\begin{itemize}
\item \textbf{Young} – 20 osób w wieku od 21 do 34 lat,
\item \textbf{Elderly} – 20 osób w wieku od 68 do 85 lat.
\end{itemize}

Każde nagranie trwa 120 minut i zostało zarejestrowane z częstotliwością próbkowania równą 250 Hz. Oprócz sygnału EKG baza zawiera adnotacje eksperckie określające położenia kolejnych załamków R, co umożliwia bezpośrednie wyznaczenie odstępów RR bez konieczności implementacji własnego algorytmu detekcji.

\subsection{Segmentacja sygnału}

Ze względu na ograniczoną liczbę dostępnych pacjentów oraz zgodnie z zaleceniami dotyczącymi krótkoterminowej analizy HRV, każde 120-minutowe nagranie podzielono na niezależne segmenty o długości 5 minut (300 s) \cite{taskforce}.
Łącznie wyekstrahowano 937 segmentów pochodzących od 40 pacjentów.

Dla każdego segmentu wyznaczano oddzielny zestaw parametrów HRV. Takie podejście pozwoliło znacząco zwiększyć liczbę próbek wykorzystywanych podczas procesu uczenia klasyfikatorów, przy jednoczesnym zachowaniu fizjologicznej interpretowalności parametrów częstotliwościowych.

\subsection{Wyznaczanie sygnału HRV}

Na podstawie czasów występowania kolejnych załamków R wyznaczono odstępy RR zgodnie z zależnością:

\[
RR_i = t_{i+1} - t_i
\]

Otrzymany szereg RR jest sygnałem nierównomiernie próbkowanym. W celu umożliwienia przeprowadzenia analizy widmowej wykonano interpolację metodą splajnów sześciennych do równomiernej siatki czasowej o częstotliwości próbkowania 4 Hz.

Następnie zastosowano procedurę usuwania trendu (\textit{detrending}), eliminującą składową stałą i wolnozmienne zmiany sygnału, które mogłyby wpływać na wyniki analizy częstotliwościowej.

\subsection{Analiza częstotliwościowa}

Gęstość widmową mocy (PSD, \textit{Power Spectral Density}) wyznaczono przy użyciu metody Welcha. Metoda ta polega na podziale sygnału na częściowo nakładające się segmenty, wyznaczeniu widma dla każdego segmentu oraz uśrednieniu otrzymanych wyników.

Na podstawie wyznaczonego PSD obliczono moc sygnału w standardowych pasmach wykorzystywanych w analizie HRV \cite{taskforce}:

\[
LF \in [0.04,\;0.15]\;Hz
\]

\[
HF \in [0.15,\;0.40]\;Hz
\]

Moc w poszczególnych pasmach wyznaczono metodą całkowania numerycznego wykorzystując funkcję \texttt{trapz}. Ze względu na duży zakres zmian wartości parametrów zastosowano transformację logarytmiczną:

\[
x_1 = \ln(LF)
\]

\[
x_2 = \ln(HF)
\]

Uzyskane wielkości wykorzystano jako cechy wejściowe klasyfikatorów.

\subsection{Klasyfikacja}

Do rozróżniania grup wiekowych zastosowano dwa klasyfikatory:

\begin{itemize}
\item \textbf{Support Vector Machine (SVM)} z jądrem radialnym RBF,
\item \textbf{Linear Discriminant Analysis (LDA)}.
\end{itemize}

Przed procesem uczenia przeprowadzono standaryzację cech metodą Z-score. Parametry normalizacji wyznaczano wyłącznie na zbiorze uczącym, a następnie wykorzystywano do przekształcenia danych testowych.

\subsection{Walidacja modeli}

Oceny skuteczności klasyfikatorów dokonano z wykorzystaniem pięciokrotnej walidacji krzyżowej typu Group K-Fold.

W odróżnieniu od klasycznej walidacji krzyżowej, wszystkie segmenty pochodzące od jednego pacjenta były zawsze przypisywane wyłącznie do zbioru uczącego lub wyłącznie do zbioru testowego. Pozwoliło to uniknąć zjawiska przecieku informacji pomiędzy zbiorami (\textit{data leakage}) i zapewniło bardziej wiarygodną ocenę zdolności uogólniania modeli.

Jako miary jakości klasyfikacji wykorzystano dokładność klasyfikacji (\textit{accuracy}) oraz macierze pomyłek (\textit{confusion matrices}).

% =========================================================

\section{Porównanie metod estymacji PSD: FFT vs Welch}
\label{sec}

Wyznaczenie parametrów częstotliwościowych HRV wymaga estymacji gęstości widmowej mocy (PSD, \textit{Power Spectral Density}) sygnału. Najczęściej stosowane są dwie metody: klasyczna analiza widmowa oparta na szybkiej transformacie Fouriera (FFT) oraz metoda Welcha. W celu wyboru najbardziej odpowiedniej techniki dla analizowanych danych przeprowadzono porównanie obu metod pod względem jakości otrzymanych widm oraz stabilności wyznaczanych parametrów LF, HF i LF/HF.

\subsection{Opis metod}

\textbf{FFT} wyznacza widmo z całego sygnału jednocześnie przy użyciu okna Hanna (w celu redukcji przecieku widmowego). Metoda charakteryzuje się wysoką rozdzielczością częstotliwościową, lecz wysoką wariancją estymatora -- widmo jest nieregularne i wrażliwe na artefakty.

\textbf{Metoda Welcha} dzieli sygnał na nakładające się segmenty, oblicza widmo każdego z nich, a następnie uśrednia wyniki. Dzięki temu wariancja estymatora jest istotnie mniejsza kosztem nieznacznego obniżenia rozdzielczości częstotliwościowej.

\subsection{Wyniki porównania}

\begin{figure}[H]
\centering
\includegraphics[width=\textwidth]{../wykresy/wykres_ekg_rpeaks.pdf}
\caption{Fragment sygnału EKG (pierwsze 60~s) z zaznaczonymi załamkami R dla rekordu \texttt{f1y01} (young) oraz \texttt{f1o01} (elderly).}
\end{figure}

\begin{figure}[H]
\centering
\includegraphics[width=\textwidth]{../wykresy/wykres_widmo_young.pdf}
\caption{Porównanie widm HRV wyznaczonych metodą FFT i Welcha dla rekordu \texttt{f1y01} (young). Kolorem niebieskim zaznaczono pasmo LF (0{,}04--0{,}15~Hz), czerwonym pasmo HF (0{,}15--0{,}40~Hz), dla pięciominutowego segmentu.}
\end{figure}

\begin{figure}[H]
\centering
\includegraphics[width=\textwidth]{../wykresy/wykres_widmo_elderly.pdf}
\caption{Porównanie widm HRV wyznaczonych metodą FFT i Welcha dla rekordu \texttt{f1o01} (elderly), dla pięciominutowego segmentu.}
\end{figure}

Wizualne rozbieżności oraz brak idealnego pokrywania się (nakładania) granic pasm niskiej częstotliwości (LF) i wysokiej częstotliwości (HF) z linią widma mocy (PSD) na wykresach wynikają z ograniczeń renderowania warstw wektorowych przez silnik graficzny środowiska MATLAB. Funkcja \texttt{area}, realizując rysowanie obiektów powierzchniowych typu wielokąt (polygon) od osi odciętych, nakłada na siebie niezależne warstwy. W procesie rasteryzacji i eksportu do formatu PDF (grafika wektorowa) powoduje to powstawanie drobnych artefaktów wizualnych na stykach zdefiniowanych zakresów częstotliwości, co jednak nie wpływa na poprawność numeryczną wyznaczonych całek (mocy pasm), na których bazowały klasyfikatory SVM i LDA.

\begin{table}[H]
\centering
\caption{Porównanie wartości parametrów HRV wyznaczonych metodą FFT i Welcha}
\begin{tabular}{|c|c|c|c|c|}
\hline
\textbf{Rekord} & \textbf{Metoda} & \textbf{LF [s²/Hz]} & \textbf{HF [s²/Hz]} & \textbf{LF/HF} \\
\hline
\texttt{f1y01} (Young) & FFT   & 0.0006 & 0.0014 & 0.4500 \\
\hline
\texttt{f1y01} (Young) & Welch & 0.0016 & 0.0035 & 0.4542 \\
\hline
\texttt{f1o01} (Elderly) & FFT   & 0.0001 & 0.0000 & 5.5170 \\
\hline
\texttt{f1o01} (Elderly) & Welch & 0.0003 & 0.0001 & 5.4619 \\
\hline
\end{tabular}
\end{table}

Mimo że wartości LF/HF uzyskane metodami FFT i Welcha są zbliżone, widma wyznaczone metodą Welcha charakteryzują się znacznie mniejszą wariancją oraz lepszą czytelnością, co ułatwia interpretację wyników i zwiększa stabilność wyznaczanych cech.

\subsection{Wybór metody}

Na podstawie przeprowadzonego porównania do dalszej analizy wybrano metodę Welch. Z następujących powodów:

\begin{itemize}
    \item Nagrania trwają 120 minut -- sygnał RR jest silnie niestacjonarny, dlatego stabilność estymatora jest ważniejsza niż rozdzielczość.
    \item Metoda Welcha daje gładkie, powtarzalne widmo, co przekłada się na stabilniejsze wartości parametrów LF oraz HF wykorzystywanych podczas procesu klasyfikacji.
    \item FFT jest wrażliwe  na artefakty ruchowe i ektopowe uderzenia serca, które mogą wystąpić podczas tak długich nagrań.
\end{itemize}

% =========================================================
\section{Wyniki badań}

W wyniku segmentacji 40 dwugodzinnych zapisów EKG wyekstrahowano łącznie 937 niezależnych segmentów o długości 5 minut. Dla każdego segmentu wyznaczono parametry LF oraz HF, które po transformacji logarytmicznej wykorzystano jako cechy wejściowe klasyfikatorów.

\subsection{Analiza widm HRV}

\begin{figure}[H]
\centering
\includegraphics[width=\textwidth]{../wykresy/widma_hrv_porownanie.pdf}
\caption{Porównanie widm HRV dla reprezentatywnych przedstawicieli grup Young i Elderly, dla pięciominutowego segmentu.}
\label{fig:widma}
\end{figure}

Na rysunku \ref{fig:widma} przedstawiono przykładowe widma HRV wyznaczone metodą Welcha dla reprezentatywnych osób należących do grup Young oraz Elderly.

Widoczne są istotne różnice w rozkładzie energii sygnału pomiędzy analizowanymi grupami. W przypadku osób młodych obserwuje się wyraźnie większą moc w paśmie HF (0.15--0.40 Hz), związanym z aktywnością układu przywspółczulnego oraz wpływem procesu oddychania na rytm serca.

Dla grupy Elderly widoczny jest spadek mocy w paśmie HF oraz względny wzrost udziału składowych niskoczęstotliwościowych LF (0.04--0.15 Hz). Wynik ten jest zgodny z obserwowanym w literaturze osłabieniem regulacji przywspółczulnej wraz z wiekiem.

\subsection{Analiza przestrzeni cech}

\begin{figure}[H]
\centering
\includegraphics[width=0.8\textwidth]{../wykresy/przestrzen_cech_scatter.pdf}
\caption{Przestrzeń cech utworzona przez parametry $\ln(LF)$ oraz $\ln(HF)$.}
\label{fig:scatter}
\end{figure}

W celu oceny separowalności klas przeanalizowano rozkład próbek w przestrzeni cech utworzonej przez logarytmy parametrów LF i HF.

Na rysunku \ref{fig:scatter} można zauważyć częściową separację pomiędzy grupami Young i Elderly. Próbki odpowiadające osobom młodym koncentrują się w obszarach wyższych wartości $\ln(HF)$, natomiast próbki należące do grupy Elderly wykazują tendencję do występowania przy niższych wartościach tej cechy.

Jednocześnie widoczne jest częściowe nakładanie się obu grup, co uzasadnia zastosowanie metod klasyfikacyjnych w celu automatycznego wyznaczenia granicy decyzyjnej pomiędzy klasami.

\subsection{Wyniki klasyfikacji}

\begin{figure}[H]
\centering
\includegraphics[width=\textwidth]{../wykresy/macierze_pomylek.pdf}
\caption{Macierze pomyłek uzyskane dla klasyfikatorów SVM oraz LDA.}
\label{fig:confusion}
\end{figure}

Klasyfikator LDA osiągnął średnią dokładność równą 78.12\%, natomiast klasyfikator SVM z jądrem RBF uzyskał dokładność 71.61\%.

Analiza macierzy pomyłek wykazała, że przewaga klasyfikatora LDA wynika przede wszystkim z lepszego rozpoznawania segmentów należących do grupy Young. Dla tej klasy uzyskano skuteczność 83.1\%, podczas gdy klasyfikator SVM osiągnął 68.9\%.

Dla grupy Elderly oba klasyfikatory uzyskały zbliżone wyniki (73.1\% dla LDA oraz 74.4\% dla SVM). Otrzymane rezultaty sugerują, że zależności pomiędzy cechami $\ln(LF)$ oraz $\ln(HF)$ a wiekiem badanych mają w dużej mierze charakter liniowy, dzięki czemu klasyfikator LDA skutecznie wyznacza granicę decyzyjną pomiędzy klasami.

\begin{table}[H]
\centering
\caption{Porównanie średniej skuteczności klasyfikatorów}
\begin{tabular}{|c|c|}
\hline
\textbf{Klasyfikator} & \textbf{Średnia dokładność [\%]} \\
\hline
SVM (RBF) & 71.61 \\
\hline
LDA & 78.12 \\
\hline
\end{tabular}
\end{table}

Uzyskane wyniki wskazują, że klasyfikator LDA osiągnął wyższą skuteczność klasyfikacji niż klasyfikator SVM. Oznacza to, że informacje zawarte w cechach $\ln(LF)$ oraz $\ln(HF)$ pozwalają na stosunkowo dobrą liniową separację grup wiekowych.

Dokładność przekraczająca 78\% potwierdza, że parametry częstotliwościowe HRV zawierają istotną informację o wieku badanych osób i mogą być wykorzystywane do budowy modeli klasyfikacyjnych.
Warto podkreślić, że ocena została przeprowadzona z wykorzystaniem walidacji Group K-Fold, co zapewnia bardziej rygorystyczną ocenę jakości modelu niż losowy podział segmentów na zbiory uczące i testowe.

% =========================================================
\section{Dyskusja wyników}

Uzyskane wyniki potwierdzają istnienie zależności pomiędzy parametrami częstotliwościowymi HRV a wiekiem badanych osób. Analiza widmowa wykazała wyraźny spadek mocy w paśmie HF u osób starszych, co jest zgodne z obserwowanym w literaturze osłabieniem aktywności układu przywspółczulnego wraz z wiekiem \cite{fantasia}.

Zastosowanie segmentacji 120-minutowych nagrań na okna pięciominutowe pozwoliło zwiększyć liczbę próbek do 937 segmentów, co umożliwiło przeprowadzenie procesu uczenia maszynowego pomimo stosunkowo niewielkiej liczby dostępnych pacjentów.

Analiza przestrzeni cech wykazała częściową separację pomiędzy grupami Young i Elderly. Oznacza to, że parametry LF oraz HF zawierają informację o wieku badanych, jednak nie pozwalają na całkowite rozdzielenie klas.

Najlepsze wyniki uzyskano dla klasyfikatora LDA, który osiągnął średnią dokładność 78.12\%. Wynik ten okazał się wyższy niż w przypadku klasyfikatora SVM z jądrem RBF. Może to sugerować, że zależności pomiędzy wykorzystanymi cechami a wiekiem mają w dużej mierze charakter liniowy.

Należy jednak zauważyć, że wykorzystano jedynie dwie cechy częstotliwościowe. Włączenie dodatkowych parametrów HRV, takich jak LF/HF, RMSSD czy SDNN, mogłoby potencjalnie zwiększyć skuteczność klasyfikacji.

% =========================================================
\section{Wnioski}

\begin{enumerate}
    \item Parametry częstotliwościowe HRV wykazują istotne różnice pomiędzy grupami Young i Elderly.
    
    \item Osoby młode charakteryzują się większą mocą w paśmie HF, związanym z aktywnością układu przywspółczulnego.
    
    \item Metoda Welcha zapewnia stabilniejszą estymację gęstości widmowej mocy niż klasyczna analiza FFT i została wybrana do dalszych analiz.
    
    \item Segmentacja sygnału na okna pięciominutowe umożliwiła zwiększenie liczby próbek do 937 segmentów przy zachowaniu zgodności z zaleceniami dotyczącymi analizy HRV.
    
    \item Najwyższą skuteczność klasyfikacji uzyskano dla klasyfikatora LDA, który osiągnął średnią dokładność 78.12\%.
    
    \item Uzyskane wyniki potwierdzają możliwość automatycznego rozróżniania grup wiekowych na podstawie parametrów częstotliwościowych HRV.
\end{enumerate}

% ========================================================

\begin{thebibliography}{9}

\bibitem{fantasia}
Iyengar N., Peng C.-K., Morin R. i in.,
Age-related alterations in the fractal scaling of cardiac interbeat interval dynamics,
American Journal of Physiology, 1996.

\bibitem{physionet}
Goldberger A.L. i in.,
PhysioBank, PhysioToolkit, and PhysioNet,
Circulation, 2000.

\bibitem{taskforce}
Task Force of the European Society of Cardiology and
the North American Society of Pacing and Electrophysiology,
Heart rate variability: Standards of measurement,
physiological interpretation and clinical use,
European Heart Journal,
vol. 17, no. 3, pp. 354--381, 1996.

\end{thebibliography}

% =========================================================

\end{document}